% Triggers both figure checks once each, and proves each is narrow.
%
% Only two lines below are findings. `step` is a pgfkeys key, so declaring a
% style with that name fails the build with an error naming the key. And the
% node text further down sets an en dash, which the prose dash check never
% sees because these files sit outside Paths.Prose. Everything else is
% ordinary TikZ that happens to mention the same words or carry the same
% characters, and none of it may be reported.
%
% Never compiled. See main.tex.
\begin{tikzpicture}[
    step/.style={draw, rounded corners=2pt},
    note/.style={font=\scriptsize, text=gray},
    thin arrow/.style={-{Stealth[length=2mm]}},
    % out/.style={dashed} would be a finding too, if it were not commented out
  ]

  \draw[step] (0,0) rectangle (2,1);
  \draw[help lines, step=0.5cm] (0,0) grid (2,1);
  \node[note, scale=0.9, text width=3cm] at (1,-1) {an ordinary node};
  \draw[thin arrow] (0,-2) -- node[midway, above] {in and out} (2,-2);

  % Path syntax, on four lines that must stay silent. `--` is how TikZ draws a
  % segment, which is the whole reason these files are outside the dash check,
  % so a node-text pass that reads paths would report every one of these.
  \draw (0,-3) -- (2,-3) -- (2,-4) -- (0,-4) -- cycle;
  \draw[thin arrow] (0,-5) -- (2,-5);
  \node at (1,-6) {a caption with an ASCII hyphen-minus in it};
  \draw (0,-7) -- (2,-7) node[right] {labelled at the end};

  % The finding: a node whose text sets an en dash.
  \node[note] at (1,-8) {parses -- then validates};

\end{tikzpicture}
